\documentclass[11pt,a4paper]{article}

% ===== Math & Symbols =====
\usepackage{amsmath,amssymb,amsthm,amsfonts}
\usepackage{mathtools}
\usepackage{mathrsfs}
\usepackage{bm}
\usepackage{bbm}

% ===== Cryptography =====
\usepackage{xstring}
\usepackage{xparse}
\usepackage[n,operators,advantage,sets,adversary,landau,probability,notions,logic,ff,mm,primitives,events,complexity,oracles,asymptotics,keys,lambda]{cryptocode}

% ===== Algorithms =====
\usepackage{algorithm}
\usepackage{algpseudocode}

% ===== Graphics & Tables =====
\usepackage{graphicx}
\usepackage{tikz}
\usetikzlibrary{arrows.meta,positioning,calc,shapes,fit}
\usepackage{booktabs}
\usepackage{array}
\usepackage{multirow}

% ===== Typography =====
\usepackage[T1]{fontenc}
\usepackage[utf8]{inputenc}
\usepackage{microtype}
\usepackage{lmodern}
\usepackage{stmaryrd}

% ===== References & Links =====
\usepackage[backend=biber,style=alphabetic]{biblatex}
\addbibresource{references.bib}
\usepackage[hidelinks]{hyperref}
\usepackage{cleveref}

% ===== Page Layout =====
\usepackage[margin=1in]{geometry}
\usepackage{enumitem}
\usepackage{xcolor}

% ===== Theorems =====
\newtheorem{theorem}{Theorem}[section]
\newtheorem{lemma}[theorem]{Lemma}
\newtheorem{proposition}[theorem]{Proposition}
\newtheorem{corollary}[theorem]{Corollary}
\theoremstyle{definition}
\newtheorem{definition}[theorem]{Definition}
\newtheorem{construction}[theorem]{Construction}
\newtheorem{assumption}[theorem]{Assumption}
\theoremstyle{remark}
\newtheorem{remark}[theorem]{Remark}
\newtheorem{example}[theorem]{Example}

% ===== Custom Macros =====
\newcommand{\F}{\mathbb{F}}
\newcommand{\Z}{\mathbb{Z}}
\newcommand{\N}{\mathbb{N}}
\newcommand{\R}{\mathbb{R}}
\newcommand{\C}{\mathbb{C}}
\newcommand{\G}{\mathbb{G}}
\newcommand{\calA}{\mathcal{A}}
\newcommand{\calB}{\mathcal{B}}
\newcommand{\calC}{\mathcal{C}}
\newcommand{\calD}{\mathcal{D}}
\newcommand{\calO}{\mathcal{O}}
\newcommand{\eps}{\varepsilon}
\providecommand{\Adv}{\mathrm{Adv}}
\providecommand{\negl}{\mathrm{negl}}
\providecommand{\poly}{\mathrm{poly}}
\newcommand{\itmac}[1]{\llbracket {#1} \rrbracket}

\providecommand{\FF}{\mathbb{F}}
\newcommand{\calP}{\mathcal{P}}
\newcommand{\calV}{\mathcal{V}}
\newcommand{\calR}{\mathcal{R}}
\newcommand{\polyrel}{\Phi}
\newcommand{\mat}[1]{\mathbf{#1}}

\title{Proving Isogeny Knowledge with VOLEitH}
\author{Hongrui Cui}
\date{\today}

\begin{document}
\maketitle

\begin{abstract}
	We describe how to use VOLEitH to prove knowledge of an isogeny
	$\phi: E_0 \to E_n$ between two supersingular elliptic curves.
	The prover commits to intermediate $j$-invariants along the isogeny
	path and proves that each consecutive pair satisfies the modular
	polynomial relation $\polyrel_{\ell}(j_{i-1}, j_i) = 0$.
\end{abstract}

\tableofcontents

% ===================================================================
\section{Introduction}
% ===================================================================

Let $E_0$ and $E_n$ be two supersingular elliptic curves over $\F_{p^2}$.
A prover $\calP$ claims to know an isogeny $\phi: E_0 \to E_n$ of degree
$\ell^k$. To prove this in zero-knowledge using VOLEitH, we decompose
the isogeny into $k$ steps of degree $\ell$ and commit to the intermediate
$j$-invariants. The verifier checks that each consecutive pair satisfies the
$\ell$-th modular polynomial $\polyrel_{\ell}$.

% ===================================================================
\section{Preliminaries}
% ===================================================================

\subsection{Modular Polynomials}

\begin{definition}[Modular Polynomial]
	For a prime $\ell$, the $\ell$-th modular polynomial
	$\polyrel_{\ell}(X, Y) \in \Z[X, Y]$ is a symmetric polynomial
	of degree $\ell + 1$ in each variable, with the property that
	\[
		\polyrel_{\ell}(j(E), j(E')) = 0
		\iff \text{there exists an isogeny } E \to E' \text{ of degree } \ell.
	\]
\end{definition}

For a small prime $\ell$, the polynomial $\polyrel_{\ell}$ can be precomputed.
For example, $\polyrel_2(X, Y)$ has the form
\begin{equation}\label{eq:modpoly}
	\polyrel_{\ell}(X, Y) = \sum_{a = 0}^{\ell+1} \sum_{b = 0}^{\ell+1}
	c_{a,b} \, X^a Y^b,
	\qquad c_{a,b} \in \Z.
\end{equation}

\subsection{Isogeny Path}

\begin{definition}[Isogeny Chain]
	An isogeny $\phi: E_0 \to E_n$ of degree $\ell^k$ can be factored as
	\[
		E_0 \xrightarrow{\phi_1} E_1 \xrightarrow{\phi_2} \cdots
		\xrightarrow{\phi_k} E_k = E_n
	\]
	where each $\phi_i: E_{i-1} \to E_i$ is an isogeny of degree $\ell$.
	Consequently, $\polyrel_{\ell}(j(E_{i-1}), j(E_i)) = 0$ for all
	$i \in [1, n]$.
\end{definition}

The prover knows the entire chain, including the intermediate curves
$E_1, \dots, E_{n-1}$. The verifier only knows $E_0$ and $E_n$.

\subsection{Batch All-but-one Vector Commitment}

A Batch All-but-one Vector Commitment (BAVC) allows a prover to commit
to $\tau$ vectors of seeds, and later open all but one position per vector.
We work over the base field $\F_{p^2}$.
The scheme is parameterized by the number of vectors $\tau$, the vector
length $N = 2^k$, and a leaf commitment function
$\mathsf{LeafCommit}: \{0,1\}^{\secpar} \to \{0,1\}^{2\secpar}$.

The BAVC provides three algorithms~\cite{faestv2}:

\begin{description}[leftmargin=0.9em, labelsep=0.5em]
	\item[$\mathsf{BAVC.Commit}(r, iv) \to (\mathsf{com}, \mathsf{decom}, \{\mathsf{sd}_{i,j}\})$:]
		On input a random root seed $r \in \{0,1\}^\lambda$ and initialization
		vector $iv$, builds a GGM tree with $L = \tau \cdot N$ leaves.
		For each leaf position $(i,j) \in [0,\tau) \times [0,N)$, computes
		$(\mathsf{sd}_{i,j}, \mathsf{com}_{i,j}) \gets \mathsf{LeafCommit}(k_\alpha, iv, \cdot)$
		where $k_\alpha$ is the leaf seed.
		Hashes the per-vector commitments $\{\mathsf{com}_{i,j}\}_j$ into
		$h_i$, then outputs the root commitment $\mathsf{com} \gets H(h_0 \| \cdots \| h_{\tau-1})$
		together with full decommitment information $\mathsf{decom}$ and all
		leaf seeds $\{\mathsf{sd}_{i,j}\}$.

	\item[$\mathsf{BAVC.Open}(\mathsf{decom}, I) \to \mathsf{decom}_I$ or $\bot$:]
		Given an index vector $I = (\Delta_0, \dots, \Delta_{\tau-1})$ specifying
		which leaf to hide in each vector, outputs a partial opening $\mathsf{decom}_I$
		that contains the unrevealed leaf commitments and exactly the GGM sibling
		seeds needed to reconstruct all other leaves.
		Aborts with $\bot$ if the required opening exceeds the threshold $T_{\sf open}$.

	\item[$\mathsf{BAVC.Reconstruct}(\mathsf{decom}_I, I, iv) \to (\mathsf{com}, \{\mathsf{sd}_{i,j}\}_{j \neq \Delta_i})$:]
		Reconstructs every leaf except the hidden ones.
		Walks up the GGM tree from the revealed sibling seeds to recover missing
		internal seeds, re-computes leaf commitments, and outputs the recomputed
		commitment $\mathsf{com}$ together with all opened leaf seeds.
\end{description}

BAVC satisfies correctness with aborts:
if $\mathsf{Open}$ returns $\mathsf{decom}_I \neq \bot$, then
$\mathsf{Reconstruct}$ recovers the original $\mathsf{com}$ and the correct
seeds $\mathsf{sd}_{i,j}$ for all $j \neq \Delta_i$.

\subsection{Information-Theoretic Message Authentication Code}

We adopt the generalized Information-Theoretic Message Authentication Code
(IT-MAC) definition of~\cite{EC:ABBS25}, adapted to our setting over~$\F_{p^2}$.

\begin{definition}\label{def:it-mac}
	For $d \in \Z$, a degree-$d$ IT-MAC for $a \in \F_{p^2}$ is
	$\itmac{a}^{d} := (p_a(x), k_a, \Delta)$,
	where a prover $\calP$ holds a polynomial ${p}_{a}(x) \in \F_{p^2}[x]$
	with $\deg(p_a) = d$ and leading coefficient $a$;
	while a verifier $\calV$ holds the evaluation
	${k}_a = {p}_{a}(\Delta)$, for a global key
	$\Delta \in \F_{p^2}$ unknown to $\calP$.
	The degree superscript $d$ is omitted when clear from context.
\end{definition}

Let $\itmac{a}^{d_a}$ and $\itmac{b}^{d_b}$ be two IT-MACs,
with corresponding polynomials $p_a(x), p_b(x)$ and evaluations $k_a, k_b$,
and assume w.l.o.g., $d_a \geq d_b$.
IT-MAC supports the following local homomorphic operations:
\begin{description}[leftmargin=0.9em, labelsep=0.5em]
	\item[Addition.] To compute $\itmac{a + b}^{d_a} = \itmac{a}^{d_a} + \itmac{b}^{d_b}$,
		$\calP$ computes $p_{a+b}(x) = p_a(x) + p_b(x) \cdot x^{d_a - d_b}$,
		while $\calV$ computes $k_{a+b} = k_a + k_b \cdot \Delta^{d_a - d_b}$.

	\item[Multiplication.] To compute $\itmac{ab}^{d_a \cdot d_b} = \itmac{a}^{d_a} \cdot \itmac{b}^{d_b}$,
		$\calP$ computes $p_{ab}(x) = p_a(x) \cdot p_b(x)$,
		while $\calV$ computes $k_{ab} = k_a \cdot k_b$.

	\item[Constants.] For any $c \in \F_{p^2}$ and $d \in \Z$,
		define $\itmac{c}^{d}$ by $p_c(x) = c x^d$ and $k_c = c \Delta^d$.
		Combining this with the above operations, constant addition and
		scalar multiplication preserve degree.
		Subtraction follows as $\itmac{a} - \itmac{b} := \itmac{a} + \itmac{-1}^{1} \cdot \itmac{b}$.
		Moreover, to de-randomize a random IT-MAC $\itmac{r}^{1}$
		(where $r \sample \F_{p^2}$) into a target value $x \in \F_{p^2}$,
		$\calP$ publishes $\delta := x - r$, and both parties compute
		$\itmac{x}^{1} := \itmac{r}^{1} + \itmac{\delta}^{1}$.
\end{description}

IT-MAC satisfies the following hiding and binding properties:
\begin{description}[leftmargin=0.9em, labelsep=0.5em]
	\item[Hiding.] Random degree-1 IT-MACs are typically generated from VOLE protocols
		(e.g., the one in \Cref{sec:voleith}), which inherently ensure
		that the committed value $a$ is independent of $\calV$'s
		authentication material. Thus, hiding is guaranteed.

	\item[Binding.] To open a committed value $a$, $\calP$ sends its
		(purported) associated polynomial $\tilde{p}_{a'}(x)$ to $\calV$,
		who checks whether $\tilde{p}_{a'}(\Delta) = k_a$.
		Any attempt by $\calP$ to open a different value $a' \neq a$
		while passing this check implies
		$\tilde{p}_{a'}(\Delta) = p_a(\Delta)$.
		Assuming $\Delta \sample \F_{p^2}$, the probability of such a
		collision is at most $\frac{\deg(p_a)}{p^2}$.
\end{description}


% ===================================================================
\section{VOLE-in-the-Head over Large Field}
\label{sec:voleith}
% ===================================================================

The VOLEitH framework has two variants, the 7-round variant specialized
for proving relations over small fields, and a 9-round variant for larger
fields~\cite{baumPubliclyVerifiableZeroKnowledge2023}.
The first variant is (somewhat) widely used in post-quantum digital
signature schemes, e.g., FAEST~v1/2, ReSolveD, and SDitH~v2 all use VOLEitH
over $\FF_{2}$.
In this write-up we focus on the 9-round variant.

\subsection{Parameters and Setup}

The prover $\calP$ holds a witness $\vec{w} \in \F^{\ell_{\sf wit}}$ that satisfies
$f_i(\vec{w}) = 0$ for all $i \in [t]$.
The verifier $\calV$ knows the constraints $\{ f_i \}_{i \in [t]}$ but not the witness.

For $n \in \mathbb{N}^+$, we define $[n] := \{ 0, 1, ..., n-1 \}$.

We define $\ell_{\sf vole} := \lceil \frac{\ell_{\sf wit} + 1}{k_{\calC}} \rceil$.
$\calP$ arranges the witness $\vec{w}$ into a $\ell_{\sf vole} \times k_{\calC}$ matrix over $\F$ as $\mat{W}$.
The internal QuickSilver proof system works on $\ell_{\sf vole} \cdot k_{\calC}$ random IT-MAC relations.
Define $\hat{\ell}_{\sf vole} := 2 \cdot \ell_{\sf vole} + 1$.
The VOLEitH framework generates $\hat{\ell}_{\sf vole}$ subspace VOLE correlations over $\F$ and we sacrifice one to check consistency.



The protocol is parameterized by:
\begin{itemize}
	\item A security parameter $\secpar$ and a field $\F = \F_{p^2}$.
	\item A pseudorandom generator $\mathsf{PRG}: \{0,1\}^{\secpar} \to \F^{\hat{\ell}_{\sf vole}}$.
	\item A number $\tau = n_C$ of parallel all-but-one VOLE instances and
		a leaf count $N = 2^k$ per instance.
	\item A systematic $[n_{\calC}, k_{\calC}, d_{\calC}]$-linear error correction code
		$\calC \subseteq \F^{n_C}$, with encoding algorithm
		$\mathsf{Enc}_{\calC}$.
	\item A small subset of $\F$
		$\calD = \{\alpha_0, \dots, \alpha_{N-1}\} \subset \F$
		of size $N = 2^k$, for converting all-but-one vector commitment into VOLE correlations.
	\item A set of degree-2 constraints $f_1, \dots, f_t$ over $\F$ to be proven. For each $i \in [1, t]$, the witness $\vec{w}$ should satisfy that $f_i (\vec{w}) = 0$.
\end{itemize}




\subsection{The 9-Round Protocol}

The public-coin protocol proceeds in nine rounds.

\begin{figure}[h]
	\centering
	\begin{tabular}{p{0.4\linewidth}p{0.2\linewidth}p{0.3\linewidth}}
		$\vec{r} \sample \{ 0,1 \}^{\secpar}$ & & \\
		$(\mathsf{com}, \mathsf{decom}, \{\mathsf{sd}_{i,j}\}) \gets \mathsf{BAVC.Com}(\vec{r})$ & & \\
	\end{tabular}
	\procedureblock{VOLEitH 9-Round Protocol}{
		\calP(\vec{w}) \< \< \calV(\{f_i\}_{i=1}^{t}) \\[][\hline]
		\vec{r} \sample \{ 0,1 \}^{\secpar} \< \< \\
		(\mathsf{com}, \mathsf{decom}, \{\mathsf{sd}_{i,j}\})
		\gets \mathsf{BAVC.Commit}(\vec{r})
		\< \< \\
		\forall i \in [\tau], \tilde{\vec{u}_{i}} := \sum_{j \in [N]} \mathsf{PRG}(\mathsf{sd}_{i,j})
		\< \< \\
	}

	\caption{The 9-round VOLEitH Protocol.}
	\label{fig:9-round-voleith-1}
\end{figure}





\begin{center}
\procedureblock[]{VOLEitH 9-Round Protocol}{%
	\calP(\text{witness } \mat{W}) \< \< \calV(\{f_i\}_{i=1}^{t}) \\[][\hline]
	r \sample \{0,1\}^{\secpar},\; iv \sample \{0,1\}^{128}
		\< \< \\
	(\mathsf{com}, \mathsf{decom}, \{\mathsf{sd}_{a,j}\})
		\gets \mathsf{BAVC.Commit}(r, iv)
		\< \< \\
	\text{For each code coordinate } a \in [n_C]\text{:} \< \< \\
	\quad (\widetilde{\vec{u}}_a,\widetilde{\vec{v}}_a)
		\gets \mathsf{SS.Sum}_{\calD}
		(\{\mathsf{sd}_{a,j}\}_{\alpha_j\in\calD})
		\< \< \\
	\widetilde{\mat{U}} \gets
		[\widetilde{\vec{u}}_1|\cdots|\widetilde{\vec{u}}_{n_C}],\quad
		\mat{V} \gets
		[\widetilde{\vec{v}}_1|\cdots|\widetilde{\vec{v}}_{n_C}]
		\< \< \\
	\mat{U} \gets \widetilde{\mat{U}}[:,1..k_C],\quad
		\mathsf{aux}_{\calC}
		\gets \mathsf{Enc}_{\calC}(\mat{U})-\widetilde{\mat{U}}
		\< \< \\
	\mat{U}=\binom{\mat{U}_1}{\mat{R}},\;
		\mat{V}=\binom{\mat{V}_1}{\mat{V}_2}
		\quad(\ell+1\text{ rows each}) \< \< \\
	\mat{D} \gets
		\mat{W} - \mat{U}_1[1..\ell]
		\< \< \\
	\< \sendmessageright*{iv,\; \mathsf{com},\; \mathsf{aux}_{\calC},\; \mat{D}} \< \\[-0.2\baselineskip]
	\chi_1,\dots,\chi_t \gets
		\mathsf{RO}_{\chi}(\mathsf{tr})
		\< \< \\
	\text{For } i \in [t]\text{, compute coefficients } \mat{A}_{i,0}, \mat{A}_{i,1} \text{ of} \< \< \\
	\quad g_i(Y) =
		\sum_{h=0}^{2} f_{i,h}(\vec{r}+\vec{w}Y)Y^{2-h}
		\< \< \\
	\mathsf{eb} \gets
		\sum_{i \in [t]}\chi_i\mat{A}_{i,0}+\vec{r}_{\ell+1}
		\< \< \\
	\mathsf{ea} \gets
		\sum_{i \in [t]}\chi_i\mat{A}_{i,1}+\vec{u}_{\ell+1}
		\< \< \\
	\< \sendmessageright*{\mathsf{eb},\; \mathsf{ea}} \< \\[-0.2\baselineskip]
	\Delta' \gets \mathsf{RO}_{\Delta'}(\mathsf{tr})
		\< \< \\
	\mat{S} \gets \mat{R} + \mat{U}_1 \cdot \Delta'
		\< \< \\
	\< \sendmessageright*{\mat{S}} \< \\[-0.2\baselineskip]
	\vec{\eta} \gets \mathsf{RO}_{\eta}(\mathsf{tr})
		\< \< \\
	\mathsf{ev} \gets
		\vec{\eta}^{\top}(\mat{V}_2+\mat{V}_1\cdot\Delta')
		\< \< \\
	\< \sendmessageright*{\mathsf{ev}} \< \\[-0.2\baselineskip]
	\vec{\delta}=(\delta_1,\dots,\delta_{n_C})
		\gets \mathsf{RO}_{\Delta}(\mathsf{tr}; [N]^{n_C})
		\< \< \\
	\vec{\Delta} \gets (\alpha_{\delta_1},\dots,\alpha_{\delta_{n_C}})
		\< \< \\
	\mathsf{decom}_{\vec{\delta}} \gets
		\mathsf{BAVC.Open}(\mathsf{decom}, \vec{\delta})
		\< \< \\
	\< \sendmessageright*{\mathsf{decom}_{\vec{\delta}}} \< \\
	\pclb\pcintertext[center]{\textbf{Verification:}}\\
	\< \<
	(\mathsf{com}', \{\mathsf{sd}_{a,j}\}_{j \neq \delta_a})
		\gets \mathsf{BAVC.Reconstruct}(\mathsf{decom}_{\vec{\delta}}, \vec{\delta}, iv) \\
	\< \<
	\text{For each } a \in [n_C]\text{: } \\
	\< \<
	\quad \widetilde{\vec{q}}_a \gets
		\mathsf{SS.Open}_{\calD}
		(\{\mathsf{sd}_{a,j}\}_{j\neq\delta_a}, \Delta_a) \\
	\< \<
	\widetilde{\mat{Q}} \gets
		[\widetilde{\vec{q}}_1|\cdots|\widetilde{\vec{q}}_{n_C}] \\
	\< \<
	\mat{Q} \gets
		\widetilde{\mat{Q}}+\mathsf{aux}_{\calC}\operatorname{diag}(\vec{\Delta}) \\
	\< \<
	\mat{Q}=\binom{\mat{Q}_1}{\mat{Q}_2},\quad
	\mat{S}' \gets \mat{S}+
		\binom{\mat{D}}{0}\cdot\Delta' \\
	\< \<
	\text{For } i\in[t]\text{, compute }
	c_i(Y)=\sum_{h=0}^{2}
		f_{i,h}(\vec{s}'_{1..\ell})Y^{2-h} \\
	\< \<
	\mathsf{es} \gets
		\sum_{i\in[t]}\chi_i c_i(\Delta')+\vec{s}'_{\ell+1} \\
	\< \<
	\text{Check } \mathsf{com}'=\mathsf{com}
		\text{ and } \mathsf{es}=\mathsf{eb}+\mathsf{ea}\cdot\Delta' \\
	\< \<
	\mat{T} \gets \mathsf{ECC.Encode}_{\calC}(\mat{S}) \\
	\< \<
	\mat{L} \gets \vec{\eta}^{\top}(\mat{Q}_2+\mat{Q}_1\cdot\Delta') \\
	\< \<
	\mat{R}_{\Delta} \gets
		\mathsf{ev}+\vec{\eta}^{\top}\mat{T}\operatorname{diag}(\vec{\Delta}) \\
	\< \<
	\text{Check } \mat{L}=\mat{R}_{\Delta} \\
	\< \< \text{Accept iff all checks pass}
}
\end{center}

\subsection{Round-by-Round Explanation}

\begin{description}[leftmargin=0.9em, labelsep=0.5em]
	\item[Commitment and subspace VOLE setup.]
		The prover commits to the $\tau=n_C$ all-but-one seed vectors with
		$\mathsf{BAVC.Commit}$.  Each vector is converted, via the
		SoftSpokenOT summation over $\calD$, into one raw coordinate of a
		VOLE transcript.  The prover takes the first $k_C$ raw coordinates
		as $\mat{U}$ and sends
		$\mathsf{aux}_{\calC}=\mathsf{Enc}_{\calC}(\mat{U})-\widetilde{\mat{U}}$
		so that the verifier can correct the opened raw VOLE values into a
		Reed--Solomon codeword.

	\item[Constraint batching.]
		The random oracle challenge $\chi$ batches the degree-2 constraints.
		The prover sends the masked QuickSilver-style coefficients
		$\mathsf{eb}$ and $\mathsf{ea}$.

	\item[Code-switch opening.]
		The random oracle challenge $\Delta'$ asks the prover to open
		$\mat{S}=\mat{R}+\mat{U}_1\Delta'$.  The additional challenge
		$\vec{\eta}$ compresses the verification of this opening to the
		single value $\mathsf{ev}$.

	\item[BAVC opening and verification.]
		The final challenge first selects indices
		$\vec{\delta}\in[N]^{n_C}$ and hence field points
		$\vec{\Delta}\in\calD^{n_C}$.  The verifier reconstructs all opened
		seeds, recomputes the raw SoftSpokenOT openings $\widetilde{\mat{Q}}$,
		corrects them as
		$\mat{Q}=\widetilde{\mat{Q}}+\mathsf{aux}_{\calC}\operatorname{diag}(\vec{\Delta})$,
		and checks both the batched constraint equation and the subspace-VOLE
		consistency equation.
\end{description}

% ===================================================================
\section{Proving Isogeny Path Knowledge}
% ===================================================================

%% TODO: Instantiate the 9-round VOLEitH framework for the isogeny path problem.
%% The witness is the chain of intermediate j-invariants j_1, ..., j_{k-1}.
%% The degree-2 constraints are the modular polynomial relations 
%% \Phi_{\ell}(j_{i-1}, j_i) = 0 for each step i \in [1, k].


% ===================================================================
\section{Security Analysis}
% ===================================================================

\begin{theorem}[Completeness]
	If $\calP$ knows a valid isogeny $\phi: E_0 \to E_n$, then the honest
	verifier accepts.
\end{theorem}

\begin{proof}
	By definition of the modular polynomial, each consecutive pair
	$(j_{i-1}, j_i)$ along the decomposed isogeny path satisfies
	$\polyrel_{\ell}(j_{i-1}, j_i) = 0$. Since VOLEitH commitments are
	computationally binding and the linear combination $s_i$ is opened
	correctly, all checks pass.
\end{proof}

\begin{theorem}[Soundness]
	If $E_0$ and $E_n$ are not connected by an isogeny of degree $\ell^k$,
	then a malicious prover $\calP^*$ cannot make $\calV$ accept except with
	negligible probability.
\end{theorem}

\begin{proof}[Proof sketch]
	If the verifier accepts, then by the binding property of VOLE
	commitments and the soundness of the linear opening, there exist
	values $j_1, \dots, j_{k-1}$ such that
	$\polyrel_{\ell}(j_{i-1}, j_i) = 0$ for all $i$.
	By the property of modular polynomials, this implies the existence
	of $\ell$-isogenies between consecutive curves, and hence a composite
	isogeny $E_0 \to E_n$ of degree $\ell^k$.
\end{proof}

\begin{remark}
	Knowledge soundness follows from the extractability of VOLEitH
	commitments: an extractor can recover $j_1, \dots, j_{k-1}$ by
	rewinding the prover and solving the resulting linear system.
\end{remark}

\printbibliography
\end{document}
